\documentclass[12pt,a4paper]{article}
\usepackage{amsmath}
\usepackage{algorithm}
\usepackage{listings}
\usepackage[noend]{algpseudocode}

\usepackage{mathpartir}
\usepackage{booktabs}
\usepackage{graphicx}
\usepackage{graphics}
\usepackage{geometry}

\usepackage{xcolor}
\usepackage{anyfontsize}

\definecolor{codegreen}{rgb}{0,0.6,0}
\definecolor{codegray}{rgb}{0.5,0.5,0.5}
\definecolor{codepurple}{rgb}{0.58,0,0.82}
\definecolor{backcolour}{rgb}{0.95,0.95,0.92}
\lstdefinestyle{mystyle}{
    backgroundcolor=\color{backcolour},
    commentstyle=\color{codegreen},
    keywordstyle=\color{magenta},
    numberstyle=\color{codegray},
    stringstyle=\color{codepurple},
    basicstyle=\ttfamily\footnotesize,
    breakatwhitespace=false,
    breaklines=true,
    captionpos=b,
    keepspaces=true,
    numbers=left,
    numbersep=5pt,
    showspaces=false,
    showstringspaces=false,
    showtabs=false,
    tabsize=2
}
\lstset{style=mystyle}

\begin{document}

\title{LLM Jailbreak Detection \& Analysis}

\author{John Hawkins}

\maketitle

\begin{abstract}
Large Large Models suffer from a vulnerability that allows malicious users
to solicit undesirable responses through manipulation of the prompt. These
so-called jailbreak prompts cause the LLM to circumvent safety guardrails put
in place to keep responses acceptable to the organisations policies.,
In this work we analyse the ability of machine learning models to recognise the
jailbreak prompts and understand what drives their behaviour.

\end{abstract}

\section{Introduction}

Large Large Models offer new opportunities for software developers to permit
natural language interfaces to multiple applications. They are being deployed
to replace search algorithms in various applications, as well as develop new 
prouctivity tools like summarisation, initial document drafting and editing. 
However, due to the nature of their training processes the models contain large 
and unquantified volumes of undersirable content sourced from across the internet. 
Explication of this undesirable content is considered
to be either socially harmful, or reputationaly damaging for the organisations
that make these models available.

Model developers are increasingly focused on preventing access to undesirable output
through either model fine tuning, or through additional detection mechanisms that
intervene before the user is provided questionable output. However, many teams have
discovered that it is possible to circumvent these safety protocols through prompt
injections that, in some sense, trick the LLM to providing content that developers
want hidden. These prompt injections have become known as jailbreaks, and a small
amount of research looks to classify and quantify the effectiveness of these prompts.

Many researchers have identified that jailbreaks tend to belong to specific semantic
families, each of which captures a particular strategy used in the prompt. The number
of strategies is increasing, as can be seen in the increase in observed categories 
provide in public datasets. In 2023 Liu \emph{et al} identifed 3 major types of
jailbreak which could be browkn down into 10 patterns\cite{Liu2023}. By 2024 
an additional X types had been observed by XXX. [Lit Review Details Here]

In this work we compile multiple jailbreak datasets and combine them with a set of
non-jailbreak prompts. Initially drawing from open source fine tuning datasets, we
extended this by crafting custom non-jailbreak prompts that in some sense mimiced the
semantic structure of the jailbreaks, by looking at legitimate uses of the general
prompt strategy. We then apply a range of machine learning technique to distinguish 
the jailbreaks from regular prompts and provide a range of insights into the features
that make certain jailbreaks difficult to detect.
 
\section{Methododology}

Split this section into two subsections, one on the dataset the other on ML approaches

\subsection{Data}

We use data from multiple existing sources, which are summarised in Table \ref{tab:sources}.

\begin{table*}
\caption{Data Sources}
\label{tab:sources}
\begin{tabular}{l|r|r}
\toprule
Source            &Jailbreak Prompts &Non-Jailbreak Prompt    \\
\midrule
Paper Authors     &73                &0              \\
\bottomrule
\end{tabular}
\end{table*}

For each of the Jailbreak datasets there are predetermined patterns of jailbreak that
the prompts are assigned to. We summarise these in Table \ref{tab:prompts}

\begin{table*}
\caption{Jailbreak Prompt Type and Patterns}
\label{tab:prompts}
\begin{tabular}{l|r|l}
\toprule
Type                   &Patterns &Description                \\
\midrule
Pretending             &3        &Pose a scenario in which the answer is innocuous.\\
Attention Shifting     &3        &Pose a scenario in which the answer is innocuous.\\
Privelidge Escalation  &3        &Pose a scenario in which the answer is innocuous.\\
\bottomrule
\end{tabular}
\end{table*}


\subsection{Models}

We test two general classes of models on this task. The first class are general machine
learning models for tabular data, using a predefined set of text features extracted
from the prompts. The second class are pre-trained BERT transformers which we then fine
tune to discriminate between jailbreak and non-jailbreak prompts.

We deliberately employ multiple pretrained BERT models of different parameter sizes to 
evaluate the impact of model size on performance.

\subsection{Evaluation}

We observe that jailbreak prompts below to different semantic classes, and 
researchers continue to identify new classes of prompt. This suggests that determining
whether it is possible for a model to detect novel classes jailbreak prompts is a 
worthwhile research question. To this end we conduct our experiments using two
forms of evaluation. In the first a simple cross validation experiment mixes all data
and evaluates the ability of each modelling approach to identify known classes
of jailbreak. 

In the second experiment, each jailbreak and related non-jailbreak are assigned to
their respective classes. We conduct training and testing such that all examples of
a particular class are either in training or in testing. The end result of which is
that test results reflect whether the model was able to infer the abstract intentional
properties of a jailbreak prompt and recignise them in novel classes of jailbreak.


\section{Results}

We break the results up into sections reflecting the different experiments

\subsection{Known Jailbreaks}

The ability of all models to identify and discriminate jailbreaks is shown in
Table \ref{results:known}

\begin{table*}
\caption{Performance on Known Jailbreak Patterns}
\label{results:known}
\begin{tabular}{l|r|r}
\toprule
Model             &AUC   &F1              \\
\midrule
Model Name        &0.6   &0.0              \\
\bottomrule
\end{tabular}
\end{table*}


\subsection{Novel Jailbreaks}

The ability of all models to identify novel jailbreak patterns, not seen during training,
 is shown in Table \ref{results:novel}

\begin{table*}
\caption{Performance on Novel Jailbreak Patterns}
\label{results:novel}
\begin{tabular}{l|r|r}
\toprule
Model             &AUC   &F1              \\
\midrule
Model Name        &0.6   &0.0              \\
\bottomrule
\end{tabular}
\end{table*}



\subsection{Jailbreak Elements}

In order to understand the linguistic elements that form a jailbreak prompt we conducted
feature analysis for these models. The results are shown below.



\section{Conclusion}

We have shown that...

\bibliographystyle{ieeetr}
\bibliography{refs}

\end{document}
